\section{Capacidad espectral, memoria energética y transporte térmico}
\label{v:sec:capacidad-espectral-termica}

La aplicación positiva entre temperatura y energía de decisión se construye
aquí sobre el portador graduado y su referencia marcada. Esta sección
conserva el origen
energético que una normalización no ve y compone esa lectura con el calendario,
el área, la gravitación y las distribuciones térmicas del artículo. La cadena
causal permanece
\[
 \mathrm{APP}\longrightarrow\mathrm{TRIT}\longrightarrow\mathrm{TPK}
 \longrightarrow X_{\rm enr}
 \longrightarrow\mathcal C_{\rm cont}^{\rm disc}
 \longrightarrow\text{portador térmico}
 \longrightarrow\text{reconocimiento físico}.
\]
Las constantes matemáticas y físicas utilizadas son salidas HMT--MD o
representaciones posteriores de ese estado. Ninguna cifra metrológica selecciona
las semillas, los sectores o el espectro. En particular, el escalar
$1{,}380649\times10^{-23}$ se obtiene primero como intensidad adimensional
marcada; su composición con la coordenatización energética y térmica se demuestra
después.

\subsection{Portador graduado y dos coordenadas conservadas}

Sean $\mathscr H_{N,L}$ cinco espacios finitos, con
\begin{equation}
 \dim\mathscr H_{0,0}=1,\quad
 \dim\mathscr H_{1,2}=380,\quad
 \dim\mathscr H_{2,0}=1,\quad
 \dim\mathscr H_{2,1}=24,\quad
 \dim\mathscr H_{2,2}=624.
 \label{v:eq:fibras-termicas-NL}
\end{equation}
Su suma directa $\mathscr H_{\rm th}$ conserva simultáneamente el grado
$N$ del prefijo y el grado $L$ de memoria incidencial. El polinomio de
Hilbert bivariado es
\begin{equation}
 Z(u,v)=1+380uv^2+u^2(1+24v+624v^2).
 \label{v:eq:polinomio-graduado-termico}
\end{equation}
Por tanto,
\[
 Z(u,1)=1+380u+649u^2,\qquad
 Z(1,v)=2+24v+1004v^2.
\]
La igualdad $649=1+24+624$ describe las fibras de $L$ en el grado $N=2$;
la identidad independiente $649=1+2\cdot324$ describe la descomposición
por involución del portador. No se intercambian esos dos significados.

Las dimensiones $380$ y $649$ proceden de dos operaciones distintas. La
primera es la parte no diagonal del endomorfismo sobre el portador de
capacidad $20$; la segunda reúne la dirección escalar y dos orientaciones
del sector de $324$ rutas. Para hacer explícita esta segunda construcción,
sea
\[
 X_\times=(\mathbb Z/9\mathbb Z)^2\times\{N,E,S,O\},
 \qquad |X_\times|=324,
\]
y sea $\mathfrak e:X_\times\to\mathcal S_\times$ el emisor de seis
ventanas, con $26$ firmas. Sus fibras tienen cardinal $8$ para $24$ firmas,
$24$ para $555555$ y $108$ para $000000$. Con
$I_{24}=\{1,\ldots,24\}$, se construye el árbol microscópico
\begin{equation}
 \widetilde{\mathcal T}=\{\Omega\}\sqcup\{a_i:i\in I_{24}\}
 \sqcup\{b_{i,y}:i\in I_{24},\ y\in X_\times\},
 \qquad |\widetilde{\mathcal T}|=7801.
 \label{v:eq:arbol-microscopico-7801}
\end{equation}
La lectura $p(b_{i,y})=b_{i,\mathfrak e(y)}$, que fija $\Omega$ y los
$a_i$, tiene imagen $\mathcal T$ de cardinal
$1+24+24\cdot26=649$. Si
$m_s=|\mathfrak e^{-1}(s)|$, la medida
\begin{equation}
 \nu_{\rm mic}(\Omega)=\nu_{\rm mic}(a_i)=1,
 \qquad \nu_{\rm mic}(b_{i,y})=m_{\mathfrak e(y)}^{-1}
 \label{v:eq:medida-arbol-microscopico}
\end{equation}
da masa uno a cada fibra de $p$. Las tasas son $r>0$ en ambos sentidos
entre $\Omega$ y $a_i$, $r/m_{\mathfrak e(y)}$ de $a_i$ a $b_{i,y}$
y $r$ en el sentido inverso.

Sobre $\ell^2(\mathcal T)$ con medida de conteo y
$\ell^2(\widetilde{\mathcal T},\nu_{\rm mic})$, el levantamiento
\begin{equation}
 J_p g=g\circ p
 \label{v:eq:levantamiento-fibras-Jp}
\end{equation}
es una isometría de la clase o firma al árbol microscópico ponderado.
En efecto,
\[
 \|J_pg\|_{\nu_{\rm mic}}^2
 =\sum_{z\in\mathcal T}|g(z)|^2
   \sum_{x:p(x)=z}\nu_{\rm mic}(x)=\|g\|^2,
 \qquad J_p^*J_p=I.
\]
El adjunto es la media ponderada
\begin{equation}
 (J_p^*f)(z)=\sum_{x:p(x)=z}\nu_{\rm mic}(x)f(x),
 \label{v:eq:adjunto-fibras-Jp}
\end{equation}
no el levantamiento $J_p$: tienen dominios opuestos y uno copia una función
de clase mientras el otro promedia su fibra. El balance detallado
$\nu_{\rm mic}(a_i)r/m_{\mathfrak e(y)}
=\nu_{\rm mic}(b_{i,y})r$ y las sumas de tasas prueban, para los
generadores microscópico y agregado,
$\widetilde L J_p=J_pL_{\mathcal T}$. El balance dentro de una fibra no
determina por sí solo los pesos entre grados; éstos se fijan por las
corrientes de borde construidas más abajo.

La variable
\begin{equation}
 Q:=N-L
 \label{v:eq:carga-memoria-Q}
\end{equation}
separa dos estados que una marginal puede identificar. Por ejemplo, conocer
$Z(u,1)$ no recupera la distribución de $L$. Las cinco celdas de
\eqref{v:eq:fibras-termicas-NL} y sus multiplicidades reconstruyen el
portador conjunto; no reconstruyen coherencias que una compresión haya
eliminado. Éste es el dominio exacto de la lectura.

\subsection{Levantamiento positivo de la memoria energética}

En las cinco celdas, $Q$ toma los valores $0,-1,2,1,0$. Introduzcamos un
registro de memoria $R$ con niveles $0,1,2,3$ y la isometría
\begin{equation}
 V_{\rm mem}:\mathscr H_{N,L}\longrightarrow
 \mathscr H_{N,L}\otimes\mathbb C^4,\qquad
 V_{\rm mem}\xi=\xi\otimes e_{N-L+1}.
 \label{v:eq:isometria-memoria-positiva}
\end{equation}
El desplazamiento en una unidad convierte la carga firmada en un nivel no
negativo sin perderla. Si $R e_j=j e_j$, entonces
\begin{equation}
 V_{\rm mem}^*(L\otimes I+I\otimes R)V_{\rm mem}=N+I.
 \label{v:eq:balance-positivo-memoria}
\end{equation}
En cada celda esta identidad es simplemente
$L+(N-L+1)=N+1$, pero su escritura operatoria conserva degeneraciones,
superposiciones dentro de cada fibra y cualquier purificación sobre un
sistema auxiliar.

Para una escala $\epsilon_0>0$, la igualdad
\begin{equation}
 V_{\rm mem}^*\epsilon_0(L\otimes I+I\otimes R)V_{\rm mem}
 =\epsilon_0(N+I)
 \label{v:eq:energia-lift-memoria}
\end{equation}
demuestra conservación exacta de energía en el levantamiento. El filtro que
retiene una rama debe conservar también el marcador de éxito; de lo contrario,
la densidad condicionada oculta el coste de preparación. El levantamiento es
reversible sobre su imagen y tiene coste lógico nulo; descartar el registro
$R$ es una operación distinta, sometida al balance de Landauer de la sección
anterior.

\subsection{Espectro decimal producido por el lazo de memoria}

La descomposición incidencial $8{:}1$ determina el acoplamiento
\[
 \mathcal Q=\frac13
 \begin{pmatrix}\sqrt8 I&-I\\ I&\sqrt8 I\end{pmatrix}.
\]
El lazo orientado de retorno es
\[
 H_1=\begin{pmatrix}2&1\\1&1\end{pmatrix},\qquad
 H_1H_1^T=\begin{pmatrix}5&3\\3&2\end{pmatrix}.
\]
En la preparación gaussiana positiva declarada, la covarianza visible
normalizada es
\begin{equation}
 W=\frac{8I+H_1H_1^T}{9},\qquad
 \det W=\frac{121}{81},\qquad
 \nu_W=\sqrt{\det W}=\frac{11}{9}.
 \label{v:eq:williamson-decimal}
\end{equation}
El teorema espectral de un modo da entonces
\begin{equation}
 r=\frac{\nu_W-1}{\nu_W+1}=\frac1{10},\qquad
 \rho_r=(1-r)\sum_{j\ge0}r^j\,|j\rangle\langle j|.
 \label{v:eq:estado-geometrico-decimal}
\end{equation}
La razón decimal se obtiene de la covarianza del lazo; no se introduce como
base numérica objetivo. El símbolo $\nu_W$ es el autovalor simpléctico de
Williamson y es distinto de la medida microscópica $\nu_{\rm mic}$ de
\eqref{v:eq:medida-arbol-microscopico}. El control $H_0=I$ produce $r=0$,
de modo que el espectro depende efectivamente de la holonomía elegida.

La división euclídea $j=3n+b$, $0\le b<3$, define el unitario
\begin{equation}
 U_3|j\rangle=|n\rangle\otimes|b\rangle,\qquad
 U_3\rho_rU_3^*=\rho_{r^3}\otimes
 \frac{\operatorname{diag}(1,r,r^2)}{1+r+r^2}.
 \label{v:eq:agrupamiento-tritico-espectral}
\end{equation}
La prueba es término a término:
$(1-r)r^{3n+b}=(1-r^3)(r^3)^n r^b/(1+r+r^2)$.
Así,
\begin{equation}
 q=r^3=10^{-3},\qquad
 \sigma_r=\frac{\operatorname{diag}(100,10,1)}{111}.
 \label{v:eq:razon-milenaria-y-residuo}
\end{equation}
El residuo conserva las tres posiciones intermedias. Si $S|j\rangle=|j+1\rangle$,
su transporte es
\[
 U_3SU_3^*=I\otimes(|1\rangle\langle0|+|2\rangle\langle1|)
 +S_n\otimes|0\rangle\langle2|,
\]
por lo que el tercer avance retorna el residuo y aumenta el cociente. Retorno
de fase y retorno del estado no se confunden.

Con $H_W=\epsilon_a(N+\tfrac12)$,
\begin{equation}
 U_3H_WU_3^*
 =\epsilon_a\left(3N_n\otimes I+I\otimes R+\tfrac12I\right).
 \label{v:eq:hamiltoniano-agrupado-decimal}
\end{equation}
La separación del factor de cociente es $3\epsilon_a$, mientras la
temperatura permanece fija:
\[
 e^{-\beta_W\epsilon_a}=10^{-1},\qquad
 e^{-\beta_W3\epsilon_a}=10^{-3},\qquad
 k_BT_W=\frac{\epsilon_a}{\log10}.
\]

\subsection{Sección marcada y evaluación completa}

La renovación de~\eqref{v:eq:tupla-renovacion-capacidad} suministra el
origen $23$ y el refinamiento por tres suministra $23,26,29$. Con
$\lambda_j/\lambda_0=r^j$ y las multiplicidades del portador,
\begin{equation}
 \boxed{
 \mathcal B
 =\frac{\lambda_{23}+380\lambda_{26}+649\lambda_{29}}{\lambda_0}
 =10^{-23}\left(1+\frac{380}{1000}+
                    \frac{649}{1000^2}\right)
 =1{,}380649\times10^{-23}.}
 \label{v:eq:evaluacion-termica-completa}
\end{equation}
El residuo $b=2$ aislado seleccionaría $2,5,8,\ldots$ y no determina el
umbral $23$; éste procede del tiempo de renovación y del índice de capacidad.
La mantisa se conserva en la densidad normalizada mediante el marcador del
sector; el peso de selección conserva además $10^{-23}$. Por ello, ni la
mantisa ni el exponente se añaden desde una tabla externa.

Sea
$\mathcal D=\bigoplus_{n=0}^2\mathcal D_n$, con
$(d_0,d_1,d_2)=(1,380,649)$,
$\dim\mathcal D_n=d_n$ y $N|_{\mathcal D_n}=nI$. Definimos el operador
de capacidad y la referencia espectral
\begin{equation}
 C=23I+3N,\qquad
 \eta=\bigoplus_{n=0}^2r^{23+3n}I_{\mathcal D_n},\qquad
 b_*:=\operatorname{Tr}\eta=\mathcal B.
 \label{v:eq:referencia-termica-eta}
\end{equation}
La referencia $\eta$ conserva por separado los grados $23,26,29$ y
permanece fija cuando varía un estado de ensayo. Para
$x=e^{-\epsilon_0/\tau}$, la partición relativa y su estado son
\begin{equation}
 Z_N(x)=1+380x+649x^2,\qquad
 \rho_N(x)=\frac{x^N}{Z_N(x)},\qquad
 p_0=\frac1{Z_N(x)}.
 \label{v:eq:particion-portador-termico}
\end{equation}
El marcador $p_0$ recupera $Z_N$ aun después de normalizar. En
$q=10^{-3}$, $Z_N(q)=1{,}380649$,
$\eta=10^{-23}q^N$ y $\eta/b_*=\rho_N(q)$; el origen marcado produce
\eqref{v:eq:evaluacion-termica-completa}. La energía libre y el grado medio son
\begin{equation}
 F_N(\tau)=-\tau\log Z_N(x),\qquad
 \langle N\rangle_x=x\frac{d}{dx}\log Z_N(x)
 =\frac{380x+1298x^2}{Z_N(x)}.
 \label{v:eq:energia-libre-grado-medio}
\end{equation}

\subsection{Corrientes de borde y rigidez del origen sectorial}

La incidencia areal proporciona $18+18$ enlaces tríticos. Sean
$u=N_{90}$, $v=N_{120}$, $a=u+v$ y
\[
 F=3u+4v,\qquad H_R=\epsilon_0F,
 \quad \epsilon_0=30\Delta\tau_R>0,
 \quad \tau_R=\frac{\hbar c}{2\pi R}.
\]
La procedencia permanece en $(u,v)$: la misma área puede corresponder a
energías diferentes. En un enlace,
$T^+=|1\rangle\langle0|+|2\rangle\langle1|$ y $T^-=(T^+)^*$ satisfacen
$[n,T^\pm]=\pm T^\pm$. Para $e,f=1,\ldots,18$, definimos
\[
 J_{ef}=T^+_{120,f}T^-_{90,e},\qquad
 B_c=\left(\sum_eT^+_{90,e}\right)^3
     \left(\sum_fT^-_{120,f}\right)^2.
\]
La regla $[F,AB]=[F,A]B+A[F,B]$ prueba
\[
 [F,J_{ef}]=J_{ef},\quad[a,J_{ef}]=0,
 \qquad [F,B_c]=B_c,\quad[a,B_c]=B_c.
\]
La secuencia de ocupaciones
$(0,2)\to(3,0)\to(2,1)$ tiene energías primitivas $8,9,10$ y áreas
$2,3,3$: el mismo incremento energético conserva dos acciones areales
distinguibles.

Los coeficientes $c_j=[z^j](1+z+z^2)^{18}$ cuentan las ocupaciones de
cada canal: $c_0=1,c_1=18,c_2=171,c_3=1122$. El carácter de frontera
\[
 P(x)=(1+x^3+x^6)^{18}(1+x^4+x^8)^{18}
\]
da $g_7=324,g_8=171,g_9=1122,g_{10}=3078$. Para los rangos
$(1,380,649)$, el primer triple admisible de energías consecutivas comienza
en $8$. Los inicios menores se excluyen porque, al usar
$g_0=1$, $g_1=g_2=g_5=0$, $g_3=g_4=18$, $g_6=171$, $g_7=324$ y
$g_8=171$, cada triple anterior incumple al menos uno de los rangos.

La compresión conectada se construye escogiendo un estado $(0,2)$,
$380$ estados $(3,0)$ y $649$ vecinos distintos de éstos mediante las
$J_{ef}$. Cada ocupación de suma tres posee al menos dos predecesores de
suma dos y cada predecesor posee como máximo $18$ sucesores. Por tanto,
los $380$ estados producen al menos
$\lceil760/18\rceil=43$ predecesores distintos y
$18\cdot43=774$ vecinos $(2,1)$, de los cuales pueden escogerse $649$.
La corriente $B_c$ conecta el estado inicial con los $380$ estados
intermedios y cada vecino final conserva un antecedente seleccionado; el
grafo resultante es conectado. Así quedan construidos, y no meramente
postulados, el subespacio comprimido, sus etiquetas y sus corrientes.

La compresión de borde realiza los grados $N=0,1,2$ a energías
$8\epsilon_0,9\epsilon_0,10\epsilon_0$. Sean
$\mathcal D_n=\ker(N-nI)$,
$P_F=\mathbf1_{\{N=1\}}$,
$B_N:\mathcal D_0\to\mathcal D_1$ y
$J_N:\mathcal D_1\to\mathcal D_2$ las corrientes comprimidas no nulas;
sus adjuntos recorren las aristas en sentido descendente. Entonces
\begin{equation}
 [N,B_N]=B_N,\quad [N,J_N]=J_N,\qquad
 [P_F,B_N]=B_N,\quad [P_F,J_N]=-J_N.
 \label{v:eq:conmutadores-corrientes-termicas}
\end{equation}
Para la corrección diagonal afín
\[
 H_\delta=\epsilon_0(8I+N)+\delta P_F
\]
se sigue de los conmutadores que
\[
 [H_\delta,B_N]=(\epsilon_0+\delta)B_N,
 \qquad
 [H_\delta,J_N]=(\epsilon_0-\delta)J_N.
\]
Exigir que ambas corrientes conserven el salto de borde $\epsilon_0$, o
simplemente que ambos saltos sean iguales, fuerza
\begin{equation}
 \delta=0.
 \label{v:eq:rigidez-desfase-corrientes}
\end{equation}
La conclusión clasifica este desfase sectorial; no afirma que se haya
clasificado todo el conmutante del Hamiltoniano.

Asignando a cada arista ascendente tasa $xr_e$ y a la descendente tasa
$r_e>0$, con el mismo $r_e$ en ambos sentidos, se obtiene balance detallado
respecto de los pesos $x^N$. La conectividad hace única la distribución
estacionaria en el portador comprimido y recupera
\eqref{v:eq:particion-portador-termico}. Esta dinámica es un semigrupo de
Markov de preparación térmica, no una evolución unitaria aislada.

\subsection{Equivalencia térmica con origen y marcador}

Sean $(\mathscr H_i,H_i,\tau_i,e_i,P_i)$ dos realizaciones, con
$\tau_i>0$, origen $e_i$, proyector de soporte $P_i$ y un isomorfismo
isométrico $J:P_1\mathscr H_1\to P_2\mathscr H_2$. La equivalencia térmica
relevante es
\begin{equation}
 J\frac{H_1-e_1I}{\tau_1}
 =\frac{H_2-e_2I}{\tau_2}J,\qquad
 JP_1=P_2J.
 \label{v:eq:equivalencia-termica-centrada}
\end{equation}
No exige $H_1=H_2$ ni $\tau_1=\tau_2$. Implica
$Jf((H_1-e_1I)/\tau_1)=f((H_2-e_2I)/\tau_2)J$ para toda función espectral
acotada $f$, y transporta los estados de Gibbs centrados. Para observables
que dependen de energía absoluta, deben transportarse además $e_i$ y la
escala: normalizar no autoriza a borrarlos.

En la compresión anterior existe una isometría basada $\iota_N$ tal que
\[
 H_R\iota_N=\iota_N\epsilon_0(8I+N).
\]
La corriente $J_N:\mathcal D_1\to\mathcal D_2$ y la isometría
$\iota_N$ son operadores distintos. Para la lectura cronológica se fijan
\begin{equation}
 C=23I+3N,\qquad
 H_{\rm cap}=\frac{\epsilon_a\log10}{\log3}C,
 \qquad \tau_{\rm clk}=\frac{\epsilon_a}{\log3}.
 \label{v:eq:lector-capacidad-escalas}
\end{equation}
La sección relativa de Williamson usa, en cambio,
$H_{W,{\rm rel}}=\epsilon_a C$ y
$\tau_W=\epsilon_a/\log10$. Por tanto,
$H_{\rm cap}/\tau_{\rm clk}=H_{W,{\rm rel}}/\tau_W=C\log10$,
aunque sus escalas de energía y temperatura son diferentes.

Con $\tau_*=\epsilon_0/\log1000$, la comparación centrada es
\begin{equation}
 \frac{H_R-8\epsilon_0I}{\tau_*}\iota_N
 =\iota_N\frac{H_{\rm cap}-23\tau_{\rm clk}\log10\,I}
                    {\tau_{\rm clk}}
 =\iota_N(\log1000)N.
 \label{v:eq:transporte-termico-centrado}
\end{equation}
La igualdad conserva grados, energías, marcador y origen sin identificar
$\epsilon_0$ con $\epsilon_a$: la primera escala pertenece a la compresión
de borde y la segunda a la sección cronológica. El factor $1/10$ entre
trazas no centradas y el desplazamiento $\log10$ son datos de la sección,
no defectos de la densidad.

\subsection{Área, procedencia y coste térmico}

En los dos canales areales, sean $N_{90},N_{120}\ge0$ y
\begin{equation}
 a=N_{90}+N_{120},\qquad
 D=90N_{90}+120N_{120}.
 \label{v:eq:ocupacion-procedencia-90-120}
\end{equation}
Se recuperan exactamente
\begin{equation}
 N_{90}=4a-\frac D{30},\qquad
 N_{120}=\frac D{30}-3a.
 \label{v:eq:inversion-procedencia-90-120}
\end{equation}
El área puede leer $a$, mientras el Hamiltoniano sectorial conserva la
procedencia:
\begin{equation}
 H_R=\tau_R\Delta D
 =\tau_R\Delta(90N_{90}+120N_{120}).
 \label{v:eq:hamiltoniano-area-procedencia}
\end{equation}
Dos estados con igual $a$ y diferente distribución entre canales tienen la
misma lectura de ocupación total y distinto coste. La reducción a un canal,
su purificación y la identidad de energía libre finita se aplican al estado
reducido ya construido; la entropía relativa cuantifica el déficit respecto
del estado de Gibbs de referencia. Esta composición explica por qué el
origen y la procedencia permanecen necesarios al relacionar área,
entrelazamiento y el funcional de Barbero--Immirzi.

\subsection{Composición con acción, gravitación y electrón}

La longitud $L_\alpha$ y la acción orientada $\hbar_a$ fijan
\begin{equation}
 \epsilon_a=\frac{\hbar_ac}{L_\alpha},\qquad
 G_a=\frac{c^3L_\alpha^2}{\hbar_a},\qquad
 k_B\Theta_{{\rm clk},a}\log3=\epsilon_a.
 \label{v:eq:triple-accion-gravedad-temperatura}
\end{equation}
Multiplicar las tres identidades elimina la sección de acción y da
\begin{equation}
 \boxed{G_a k_B\Theta_{{\rm clk},a}\log3=c^4L_\alpha.}
 \label{v:eq:identidad-gravedad-temperatura}
\end{equation}
La identidad se obtiene en la misma orientación $a$; $G_a$ no se introduce
desde metrología. La letra $L_\alpha$ designa longitud y no el grado $L$ de
\eqref{v:eq:fibras-termicas-NL}.

Si $R_{\rm act}=\hbar_{\rm pre}/\hbar_{\rm ret}$ es la razón de acción del
retorno electrónico, la misma composición produce
\begin{equation}
 R_{\rm act}
 =\frac{\Theta_{{\rm clk,pre}}}{\Theta_{{\rm clk,ret}}}
 =\frac{G_{\rm ret}}{G_{\rm pre}}.
 \label{v:eq:retorno-accion-termico-gravitatorio}
\end{equation}
Esta igualdad transporta la sección térmica; no usa $k_B$ para seleccionar
la ruta ni rederiva en este artículo la ley de masas. Para una energía
electrónica $E_e$ ya individualizada, $y_e=E_e/\epsilon_{\rm ret}$ satisface
\[
 \beta_{\rm clk}E_e=(\log3)y_e,\qquad
 \frac{G_{\rm ret}m_e^2}{\hbar_{\rm ret}c}=y_e^2.
\]
El origen de energía afecta masa y radio aunque se cancele en probabilidades
normalizadas.

\subsection{Puente termométrico marcado e identidad efectiva}

La referencia $\eta$ de~\eqref{v:eq:referencia-termica-eta} debe permanecer
separada del estado sobre el que se mide energía. Sea $\mathcal H$ un espacio
finito, $\rho$ una matriz de densidad sobre $\mathcal H$ y
$H=H^*$ un Hamiltoniano con origen registrado. En
$\widehat{\mathcal D}=\mathcal D\oplus\mathbb C f$ definimos
\begin{equation}
 \widehat\eta=\eta\oplus(1-b_*)|f\rangle\langle f|,
 \qquad \Omega_\rho=\widehat\eta\otimes\rho,
 \label{v:eq:referencia-marcada-normalizada}
\end{equation}
y denotamos por $P$ el proyector sobre el sumando $\mathcal D$, extendido
por la identidad de $\mathcal H$. Con
$s(\rho)=-\operatorname{Tr}(\rho\log\rho)$ y los logaritmos restringidos
al soporte, las dos lecturas locales son
\begin{align}
 \mathcal S_P(\rho)
 &=-\operatorname{Tr}\!\left[P\Omega_\rho
                 (I\otimes\log\rho)\right]
   =b_*s(\rho),\label{v:eq:lector-entropico-marcado}\\
 \mathcal E_P(\rho)
 &=\frac{\operatorname{Tr}\!\left[P\Omega_\rho(I\otimes H)\right]}
         {\operatorname{Tr}(P\Omega_\rho)}
   =\operatorname{Tr}(\rho H).
 \label{v:eq:lector-energetico-condicionado}
\end{align}
En efecto, $P\Omega_\rho=\eta\otimes\rho$ y
$\operatorname{Tr}(P\Omega_\rho)=b_*$. La factorización de la traza da
$b_*s(\rho)$ y $b_*\operatorname{Tr}(\rho H)$ en los dos numeradores;
dividir el segundo por $b_*>0$ prueba la lectura energética. Así se
conservan simultáneamente la referencia fija, la información relativa a
su marca y la energía por estado condicionado.

Trabajamos en las bases positivas transportadas de las rectas de energía,
temperatura y entropía, con $u_{\rm ent}=u_E/u_\Theta$. En las fórmulas de
coordenadas, $H$ expresa energía en $u_E$ y $\mathcal S_P$ expresa entropía
en $u_{\rm ent}$. La regla de realización consta de cuatro operaciones:
recibir la referencia fija~\eqref{v:eq:referencia-marcada-normalizada};
evaluar~\eqref{v:eq:lector-entropico-marcado};
evaluar~\eqref{v:eq:lector-energetico-condicionado}; y definir la sección
térmica como conjugada de esos dos funcionales. Las bases se transportan con
los operadores, sin introducir una temperatura objetivo.

\begin{theorem}[Derivada termométrica de la lectura marcada]
\label{v:thm:derivada-termometrica-marcada}
Si $H$ no es escalar y
$\rho_\beta=e^{-\beta H}/Z(\beta)$ para $\beta>0$, entonces
\begin{equation}
 \frac{d\mathcal S_P}{d\mathcal E_P}=b_*\beta,
 \qquad \Theta_P=\frac1{b_*\beta},
 \qquad B_*(\Theta_Pu_\Theta)=\beta^{-1}u_E,
 \label{v:eq:derivada-termometrica-marcada}
\end{equation}
donde $B_*(u_\Theta)=b_*u_E$ es la única aplicación lineal positiva con esa
imagen de la base.
\end{theorem}
\begin{proof}
Para $E(\beta)=\operatorname{Tr}(\rho_\beta H)$, la suma espectral finita
da $E'=-\operatorname{Var}_{\rho_\beta}(H)<0$ y
$(\log Z)'=-E$. De
$s(\rho_\beta)=\beta E+\log Z$ resulta $s'=\beta E'$. Como $\eta$ se
mantiene fija,
$\mathcal S_P'=b_*s'$ y $\mathcal E_P'=E'$. Su cociente y su inverso
prueban las dos primeras identidades; la imagen de una base positiva fija
una única aplicación lineal y prueba la tercera.
\end{proof}

El factor depende de la pareja de lecturas. Si energía y entropía se toman
ambas marcadas, $(b_*E,b_*s)$, o ambas condicionadas, $(E,s)$, la derivada
es $\beta$; para la pareja marcada--condicionada
\eqref{v:eq:lector-entropico-marcado}--
\eqref{v:eq:lector-energetico-condicionado} es $b_*\beta$.

La misma regla posee una caracterización variacional. Para $\Theta>0$,
\begin{equation}
 \Phi_\Theta(\rho)=\mathcal E_P(\rho)-\Theta\mathcal S_P(\rho)
 =\operatorname{Tr}(\rho H)-b_*\Theta s(\rho)
 \label{v:eq:funcional-variacional-marcado}
\end{equation}
tiene el único minimizador
\begin{equation}
 \rho_\Theta^*=\frac{e^{-H/(b_*\Theta)}}
                    {\operatorname{Tr}e^{-H/(b_*\Theta)}}.
 \label{v:eq:minimo-variacional-marcado}
\end{equation}
En efecto,
\[
 \Phi_\Theta(\rho)-\Phi_\Theta(\rho_\Theta^*)
 =b_*\Theta D_{\rm rel}(\rho\Vert\rho_\Theta^*)\ge0,
\]
con
$D_{\rm rel}(\rho\Vert\sigma)
=\operatorname{Tr}[\rho(\log\rho-\log\sigma)]$; la igualdad en la
positividad de la entropía relativa ocurre exactamente
para $\rho=\rho_\Theta^*$. Esto demuestra existencia y unicidad sin
introducir una temperatura objetivo.

Aplicamos ahora el resultado al reloj de capacidad. Fijadas la orientación
y $\epsilon_a$, sean $H_{{\rm clk},a}=H_{\rm cap}$ y
$\beta_{{\rm clk},a}=\log3/\epsilon_a$. Este Hamiltoniano es positivo,
autoadjunto y no escalar, pues sus tres autovalores distintos son
proporcionales a $23,26,29$. Entonces
\begin{align}
 \beta_{{\rm clk},a}H_{{\rm clk},a}
   &=C\log10,\nonumber\\
 e^{-\beta_{{\rm clk},a}H_{{\rm clk},a}}
   &=10^{-23}q^N=\eta,\qquad q=10^{-3},\nonumber\\
 \rho_{{\rm clk},a}
   &=\frac{q^N}{Z_N(q)}=\frac{\eta}{b_*}.
 \label{v:eq:estado-reloj-capacidad}
\end{align}
La referencia $\eta$, independiente de $\epsilon_a$, permanece fija en el
primer factor, mientras $\rho_\beta$ varía en una segunda copia de
$\mathcal D$ con $H_{{\rm clk},a}$ fijo. Evaluar después en
$\beta=\beta_{{\rm clk},a}$ da
\begin{equation}
 \Theta_{{\rm clk},P,a}=\frac{\epsilon_a}{b_*\log3},
 \qquad
 \boxed{B_*(\Theta_{{\rm clk},P,a}u_\Theta)
 =\frac{\epsilon_a}{\log3}u_E
 =\frac{h_a}{108t_0\log3}u_E.}
 \label{v:eq:identidad-transductor-numeral}
\end{equation}
La última igualdad usa la coordenada de acción ya fijada
$\epsilon_a=h_a/(108t_0)$; no usa el valor de $b_*$ como ajuste. Tampoco
identifica $\epsilon_a$ con la escala de borde $\epsilon_0$.

Si $B_{\rm clk}$ es el transductor cronológico y $J_E,J_\Theta$ transportan
respectivamente energía y las dos secciones termométricas del mismo estado
de referencia, la unicidad entre rectas da el cuadrado tipado
\begin{equation}
 B_*J_\Theta=J_EB_{\rm clk}.
 \label{v:eq:cuadrado-transductor-marcado}
\end{equation}
Las dos composiciones coinciden en la sección positiva generadora por
\eqref{v:eq:identidad-transductor-numeral}; la linealidad extiende la
igualdad a toda la recta. Éste es el puente entre la evaluación espectral
de~\eqref{v:eq:evaluacion-termica-completa} y el transductor dimensional:
dominios, referencia, lectores y sección positiva determinan el mapa antes
de su evaluación numérica.

\subsection{Consecuencias en ocupación, radiación e información}

El portador finito aporta a la teoría de Gibbs la partición
$Z_N(x)$, el estado $\rho_N(x)$, el marcador $p_0$ y la energía libre
\eqref{v:eq:energia-libre-grado-medio}. Las completaciones simétrica y
exterior desarrolladas más adelante actúan sobre modos completos y producen
las leyes de Bose--Einstein y Fermi--Dirac; el portador $1/380/649$ no las
sustituye. Les suministra una sección energética y una preparación con origen
registrado. En el régimen diluido, la misma coordenada
$x=e^{-\beta(E-\mu)}$ entra en el límite Maxwell--Boltzmann ya demostrado.

En la realización radiativa, la composición
$\beta=(k_B\Theta)^{-1}$ y
$k_B\Theta_{{\rm clk},a}\log3=\epsilon_a$ sustituye la escala del calendario
en la densidad de Planck sin alterar el conteo de dos polarizaciones ni su
dispersión. Los momentos conservan las deducciones de Stefan--Boltzmann y
Wien; la nueva sección determina de dónde procede la escala térmica y qué
marcador debe conservarse al condicionar una rama. En el área y el
entrelazamiento,~\eqref{v:eq:inversion-procedencia-90-120} impide que la
ocupación total borre el coste de procedencia. Así, ocupación, radiación,
energía libre, área y gravitación reciben la misma información térmica sin
confundir sus observables.
